\begin{PSbox}[unbreakable]{obj}{Learning Objectives}

After completing this module, students will be able to:

\begin{enumerate}
\def\labelenumi{\arabic{enumi}.}
\tightlist
\item
  Identify which probability distribution models a given engineering scenario
\item
  State the $\frac{\mathrm{PMF}}{\mathrm{PDF}}$, CDF, mean, and variance for all major distributions
\item
  Explain the physical interpretation and engineering context for each distribution
\item
  Determine when to use (and when NOT to use) each distribution
\item
  Analyze and compare probability distributions
\item
  Recognize distribution patterns from data characteristics
\end{enumerate}

\end{PSbox}

\PSsection{3.1}{Introduction: Choosing the Right Distribution}

\PSsubsection{The Engineering Challenge}

When modeling a random phenomenon, the critical question is:

\begin{PSquote}

``Which probability distribution best describes this uncertain quantity?''

\end{PSquote}

Choosing the wrong distribution leads to incorrect predictions, improper designs, and system failures.

\PSsubsection{How to Choose}

The distribution should match the \textbf{physical mechanism} generating randomness:

\begin{itemize}
\tightlist
\item
  Counting successes \ensuremath{\to} Binomial
\item
  Waiting for events \ensuremath{\to} Exponential/Geometric
\item
  Sum of many small effects \ensuremath{\to} Gaussian
\item
  Signal envelope in fading \ensuremath{\to} Rayleigh
\end{itemize}

This module provides the tools to make this choice correctly.

\PSsection{3.2}{Bernoulli Distribution}

\PSsubsection{Physical Interpretation}

Models a \textbf{single trial} with two possible outcomes: ``success'' (1) or ``failure'' (0).

\PSsubsection{Engineering Context}

Any binary random outcome:

\begin{itemize}
\tightlist
\item
  One bit transmitted: error or correct
\item
  One component tested: pass or fail
\item
  One packet sent: delivered or lost
\item
  One detection attempt: target present or absent
\end{itemize}

\begin{PSbox}[unbreakable]{def}{Mathematical Definition}

$X \sim \mathrm{Bernoulli}(p)$, where $p$ = probability of success.

\end{PSbox}

\PSsubsection{PMF}

\PSdisplay{p_X(x) = \begin{cases} p & x = 1 \\ 1-p = q & x = 0 \end{cases}}

Compact form: $p_{X}(x) = p^{x}(1-p)^{1-x}$ for $x \in \{0,\, 1\}$

\PSsubsection{CDF}

\PSdisplay{F_X(x) = \begin{cases} 0 & x < 0 \\ 1-p & 0 \leq x < 1 \\ 1 & x \geq 1 \end{cases}}

\PSsubsection{Mean}

\PSdisplay{E[X] = p}

\PSsubsection{Variance}

\PSdisplay{\text{Var}(X) = p(1-p)}

Maximum variance at $p = 0.5$ (maximum uncertainty).

\PSsubsection{When to Use}

\PSyes{} Single binary trial with fixed probability\PSbr{}\PSyes{} Building block for more complex distributions\PSbr{}\PSyes{} Modeling $\frac{ON}{\mathrm{OFF}}$ states, pass/fail outcomes

\PSsubsection{When NOT to Use}

\PSno{} Multiple trials (use Binomial)\PSbr{}\PSno{} Outcome is not binary\PSbr{}\PSno{} Probability changes between trials

\PSsection{3.3}{Binomial Distribution}

\PSsubsection{Physical Interpretation}

Models the \textbf{number of successes} in $n$ \textbf{independent} trials, each with the same success probability $p$.

\PSsubsection{Engineering Context}

\begin{itemize}
\tightlist
\item
  Number of bit errors in $n$ transmitted bits
\item
  Number of defective components in a batch of $n$
\item
  Number of successful packet transmissions out of $n$ attempts
\item
  Number of sensors that detect a signal (out of $n$ sensors)
\end{itemize}

\begin{PSbox}[unbreakable]{def}{Mathematical Definition}

$X \sim \mathrm{Binomial}(n,\, p)$

\end{PSbox}

\PSsubsection{PMF}

\PSdisplay{p_X(k) = \binom{n}{k} p^k (1-p)^{n-k}, \quad k = 0, 1, 2, \ldots, n}

where $C(n,\,k) = \frac{n!}{k!(n-k)!}$

\PSsubsection{CDF}

\PSdisplay{F_X(k) = \sum_{i=0}^{\lfloor k \rfloor} \binom{n}{i} p^i (1-p)^{n-i}}

No closed-form expression — use tables or numerical summation.

\PSsubsection{Mean}

\PSdisplay{E[X] = np}

\PSsubsection{Variance}

\PSdisplay{\text{Var}(X) = np(1-p)}

\PSsubsection{When to Use}

\PSyes{} Fixed number of trials $n$\PSbr{}\PSyes{} Each trial is independent\PSbr{}\PSyes{} Each trial has the same probability $p$\PSbr{}\PSyes{} Counting the number of ``successes''

\PSsubsection{When NOT to Use}

\PSno{} Trials are not independent (correlated failures)\PSbr{}\PSno{} Probability changes between trials\PSbr{}\PSno{} $n$ is not fixed (use Negative Binomial or Poisson)\PSbr{}\PSno{} Very large $n$ with small $p$ (use Poisson approximation)

\PSsection{3.4}{Geometric Distribution}

\PSsubsection{Physical Interpretation}

Models the \textbf{number of trials until the first success} (or equivalently, the number of failures before the first success).

\PSsubsection{Engineering Context}

\begin{itemize}
\tightlist
\item
  Number of transmissions until a packet is successfully delivered
\item
  Number of components tested until finding a defective one
\item
  Number of login attempts until successful authentication
\item
  Number of time slots until channel becomes available
\end{itemize}

\begin{PSbox}[unbreakable]{def}{Mathematical Definition}

$X \sim \mathrm{Geometric}(p)$

Convention: $X$ = number of the trial on which first success occurs $(X \in \{1,\, 2,\, 3,\, \ldots \})$.

\end{PSbox}

\PSsubsection{PMF}

\PSdisplay{p_X(k) = (1-p)^{k-1} p, \quad k = 1, 2, 3, \ldots}

($k-1$ failures, then one success)

\PSsubsection{CDF}

\PSdisplay{F_X(k) = 1 - (1-p)^k, \quad k = 1, 2, 3, \ldots}

\PSsubsection{Mean}

\PSdisplay{E[X] = \frac{1}{p}}

\PSsubsection{Variance}

\PSdisplay{\text{Var}(X) = \frac{1-p}{p^2}}

\PSsubsection{Memoryless Property}

The geometric distribution is the ONLY discrete distribution with the memoryless property:

\PSdisplay{P(X > m + n \mid X > m) = P(X > n)}

\textbf{Engineering meaning:} If a system hasn’t succeeded in $m$ attempts, the probability of needing $n$ more attempts is the same as starting fresh. ``Past failures don’t affect future probability.''

\PSsubsection{When to Use}

\PSyes{} Repeated independent trials until first success\PSbr{}\PSyes{} Each trial has the same probability\PSbr{}\PSyes{} Memoryless assumption is appropriate (e.g., independent retransmissions)

\PSsubsection{When NOT to Use}

\PSno{} Success probability changes over time (aging, learning)\PSbr{}\PSno{} Trials are not independent\PSbr{}\PSno{} Looking for the $r-\mathrm{th}$ success (use Negative Binomial)

\PSsection{3.5}{Negative Binomial Distribution}

\PSsubsection{Physical Interpretation}

Models the \textbf{number of trials until the $r-\mathrm{th}$ success} (generalization of Geometric).

\PSsubsection{Engineering Context}

\begin{itemize}
\tightlist
\item
  Number of transmissions until $r$ packets are successfully delivered
\item
  Number of components inspected until finding $r$ defectives
\item
  Number of time slots until $r$ channels become available
\end{itemize}

\begin{PSbox}[unbreakable]{def}{Mathematical Definition}

$X \sim \mathrm{NegBin}(r,\, p)$ = number of trials until $r-\mathrm{th}$ success.

\end{PSbox}

\PSsubsection{PMF}

\PSdisplay{p_X(k) = \binom{k-1}{r-1} p^r (1-p)^{k-r}, \quad k = r, r+1, r+2, \ldots}

(Choose which $r-1$ of the first $k-1$ trials were successes, the $k-\mathrm{th}$ is the $r-\mathrm{th}$ success.)

\PSsubsection{CDF}

No simple closed form. Computed by summation or numerical methods.

\PSsubsection{Mean}

\PSdisplay{E[X] = \frac{r}{p}}

\PSsubsection{Variance}

\PSdisplay{\text{Var}(X) = \frac{r(1-p)}{p^2}}

\PSsubsection{Relationship to Geometric}

When $r = 1$, the Negative Binomial reduces to the Geometric distribution.

\PSsubsection{When to Use}

\PSyes{} Repeated independent trials until $r-\mathrm{th}$ success\PSbr{}\PSyes{} Constant probability per trial\PSbr{}\PSyes{} Generalizing geometric to multiple required successes

\PSsubsection{When NOT to Use}

\PSno{} Only need first success (use Geometric — simpler)\PSbr{}\PSno{} Trials not independent\PSbr{}\PSno{} Fixed number of trials (use Binomial)

\PSsection{3.6}{Poisson Distribution}

\PSsubsection{Physical Interpretation}

Models the \textbf{number of events occurring in a fixed interval} (time, space, area) when events occur independently at a constant average rate.

\PSsubsection{Engineering Context}

\begin{itemize}
\tightlist
\item
  Number of photons hitting a detector in $1 \mu s$
\item
  Number of packets arriving at a router in 1 second
\item
  Number of defects per $cm^{2}$ on a silicon wafer
\item
  Number of cosmic ray events in a sensor per hour
\item
  Number of interference spikes in a measurement window
\end{itemize}

\begin{PSbox}[unbreakable]{def}{Mathematical Definition}

$X \sim \mathrm{Poisson}(\lambda)$, where $\lambda$ = average number of events per interval.

\end{PSbox}

\PSsubsection{PMF}

\PSdisplay{p_X(k) = \frac{\lambda^k e^{-\lambda}}{k!}, \quad k = 0, 1, 2, \ldots}

\PSsubsection{CDF}

\PSdisplay{F_X(k) = e^{-\lambda} \sum_{i=0}^{\lfloor k \rfloor} \frac{\lambda^i}{i!}}

\PSsubsection{Mean}

\PSdisplay{E[X] = \lambda}

\PSsubsection{Variance}

\PSdisplay{\text{Var}(X) = \lambda}

\textbf{Unique property:} Mean equals variance! This is a quick check — if sample mean \ensuremath{\approx} sample variance, data may be Poisson.

\PSsubsection{Poisson as Binomial Approximation}

When $n$ is large and $p$ is small, with $\lambda$ = np:

\PSdisplay{\text{Binomial}(n, p) \approx \text{Poisson}(\lambda = np)}

Rule of thumb: $n \ge 20$ and $p \le 0.05$.

\PSsubsection{When to Use}

\PSyes{} Counting events in a fixed interval\PSbr{}\PSyes{} Events occur independently\PSbr{}\PSyes{} Events occur at a constant average rate\PSbr{}\PSyes{} Rare events with many opportunities (large $n$, small $p$)\PSbr{}\PSyes{} Mean \ensuremath{\approx} Variance in data

\PSsubsection{When NOT to Use}

\PSno{} Events are not independent (clustered arrivals)\PSbr{}\PSno{} Rate varies with time (non-homogeneous process)\PSbr{}\PSno{} Events are not rare ($p$ not small) — use Binomial\PSbr{}\PSno{} Mean \ensuremath{\neq} Variance (look at overdispersion/underdispersion)

\PSsection{3.7}{Uniform Distribution (Continuous)}

\PSsubsection{Physical Interpretation}

Models a quantity that is \textbf{equally likely} to take any value in an interval $[a,\, b]$. Represents \textbf{maximum uncertainty} within a known range.

\PSsubsection{Engineering Context}

\begin{itemize}
\tightlist
\item
  Phase of a received signal (unknown, equally likely in $[0,\, 2\pi)$)
\item
  Round-off error in an ADC (uniform in $\left[-\frac{\Delta}{2},\, \frac{\Delta}{2}\right]$ where $\Delta$ = quantization step)
\item
  Arrival time within a time slot (unknown position)
\item
  Random access protocol: choose a random backoff time
\end{itemize}

\begin{PSbox}[unbreakable]{def}{Mathematical Definition}

$X \sim \mathrm{Uniform}(a,\, b)$ or $X \sim U(a,\, b)$

\end{PSbox}

\PSsubsection{PDF}

\PSdisplay{f_X(x) = \begin{cases} \frac{1}{b-a} & a \leq x \leq b \\ 0 & \text{otherwise} \end{cases}}

\PSsubsection{CDF}

\PSdisplay{F_X(x) = \begin{cases} 0 & x < a \\ \frac{x-a}{b-a} & a \leq x \leq b \\ 1 & x > b \end{cases}}

\PSsubsection{Mean}

\PSdisplay{E[X] = \frac{a+b}{2}}

\PSsubsection{Variance}

\PSdisplay{\text{Var}(X) = \frac{(b-a)^2}{12}}

\PSsubsection{When to Use}

\PSyes{} All values in a range are equally likely\PSbr{}\PSyes{} No information favors any particular value\PSbr{}\PSyes{} Quantization error modeling\PSbr{}\PSyes{} Phase of unmodulated carrier

\PSsubsection{When NOT to Use}

\PSno{} Values near the center are more likely (use Gaussian)\PSbr{}\PSno{} Values are non-negative with exponential decay (use Exponential)\PSbr{}\PSno{} Range is unbounded

\PSsection{3.8}{Exponential Distribution}

\PSsubsection{Physical Interpretation}

Models the \textbf{time between events} in a Poisson process, or the \textbf{lifetime} of a memoryless component.

\PSsubsection{Engineering Context}

\begin{itemize}
\tightlist
\item
  Time between packet arrivals
\item
  Time until component failure (constant failure rate)
\item
  Time between cosmic ray events in a detector
\item
  Inter-arrival time of calls at a switch
\item
  Duration of a fading dip in a wireless channel
\end{itemize}

\begin{PSbox}[unbreakable]{def}{Mathematical Definition}

$X \sim \mathrm{Exponential}(\lambda)$ where $\lambda$ = rate parameter (events per unit time).

Alternative parameterization: $X \sim \mathrm{Exp}(\beta)$ where $\beta = \frac{1}{\lambda}$ = mean.

\end{PSbox}

\PSsubsection{PDF}

\PSdisplay{f_X(x) = \lambda e^{-\lambda x}, \quad x \geq 0}

Or with mean $\beta$: $f_{X}(x) = \frac{1}{\beta}e^{-x/\beta}$

\PSsubsection{CDF}

\PSdisplay{F_X(x) = 1 - e^{-\lambda x}, \quad x \geq 0}

\PSsubsection{Mean}

\PSdisplay{E[X] = \frac{1}{\lambda} = \beta}

\PSsubsection{Variance}

\PSdisplay{\text{Var}(X) = \frac{1}{\lambda^2} = \beta^2}

\textbf{Note:} Standard deviation = mean for exponential distribution.

\PSsubsection{Memoryless Property}

\PSdisplay{P(X > s + t \mid X > s) = P(X > t)}

The exponential is the ONLY continuous distribution with this property.

\textbf{Engineering meaning:} A component that has been working for $s$ hours is statistically ``as good as new.'' Its remaining lifetime has the same distribution as a brand-new component.

\PSsubsection{Relationship to Poisson}

If events arrive as a Poisson process with rate $\lambda$:

\begin{itemize}
\tightlist
\item
  Number of events in time $T \sim \mathrm{Poisson}(\lambda T)$
\item
  Time between events $\sim \mathrm{Exponential}(\lambda)$
\end{itemize}

\PSsubsection{When to Use}

\PSyes{} Modeling time between events in a Poisson process\PSbr{}\PSyes{} Component lifetime with constant failure rate\PSbr{}\PSyes{} Memoryless assumption is reasonable

\PSsubsection{When NOT to Use}

\PSno{} Failure rate increases with age (wear-out) — use Weibull or Gamma\PSbr{}\PSno{} Failure rate decreases with age (burn-in) — use Weibull\PSbr{}\PSno{} Data shows increasing hazard — NOT memoryless

\PSsection{3.9}{Gaussian (Normal) Distribution}

\PSsubsection{Physical Interpretation}

Models quantities that result from the \textbf{sum of many small, independent random effects}. The most important distribution in engineering due to the Central Limit Theorem.

\PSsubsection{Engineering Context}

\begin{itemize}
\tightlist
\item
  Thermal noise voltage (sum of many electron contributions)
\item
  Measurement errors (sum of many small error sources)
\item
  Manufacturing variations (many independent process variations)
\item
  Aggregate interference in communications
\item
  Any quantity well-modeled by CLT
\end{itemize}

\begin{PSbox}[unbreakable]{def}{Mathematical Definition}

$X \sim N(\mu,\, \sigma^{2})$ where $\mu$ = mean, $\sigma^{2}$ = variance.

\end{PSbox}

\PSsubsection{PDF}

\PSdisplay{f_X(x) = \frac{1}{\sigma\sqrt{2\pi}} \exp\left(-\frac{(x-\mu)^2}{2\sigma^2}\right), \quad -\infty < x < \infty}

\PSsubsection{CDF}

\PSdisplay{F_X(x) = \Phi\left(\frac{x-\mu}{\sigma}\right) = \frac{1}{2}\left[1 + \text{erf}\left(\frac{x-\mu}{\sigma\sqrt{2}}\right)\right]}

No closed-form — computed numerically or via tables.

\PSsubsection{Mean}

\PSdisplay{E[X] = \mu}

\PSsubsection{Variance}

\PSdisplay{\text{Var}(X) = \sigma^2}

\PSsubsection{Standard Normal}

$Z \sim N(0,\, 1)$: standardized form.

Any Normal can be standardized: $Z = \frac{X - \mu}{\sigma}$

\PSsubsection{Key Probability Rules $(\sigma -\text{rules})$}

{\def\LTcaptype{none} % do not increment counter
\begin{longtable}[c]{@{}cc@{}}
\toprule\noalign{}
\rowcolor{PSnavy}\PSth{Range} & \PSth{Probability} \\
\midrule\noalign{}
\endhead
\bottomrule\noalign{}
\endlastfoot
$\mu \pm 1\sigma$ & 68.27\% \\
$\mu \pm 2\sigma$ & 95.45\% \\
$\mu \pm 3\sigma$ & 99.73\% \\
$\mu \pm 4\sigma$ & 99.9937\% \\
\end{longtable}
}

\PSsubsection{Properties}

\begin{enumerate}
\def\labelenumi{\arabic{enumi}.}
\tightlist
\item
  \textbf{Symmetric} about $\mu$
\item
  \textbf{Linear transformation:} If $X \sim N(\mu,\, \sigma^{2})$, then aX $+ b \sim N(a\mu + b,\, a^{2}\sigma^{2})$
\item
  \textbf{Sum of Gaussians:} If $X \sim N(\mu_{1},\, \sigma_{1}^{2})$ and $Y \sim N(\mu_{2},\, \sigma_{2}^{2})$ are independent, then $X + Y \sim N(\mu_{1}+\mu_{2},\, \sigma_{1}^{2}+\sigma_{2}^{2})$
\item
  \textbf{Fully characterized} by mean and variance (all higher moments determined)
\end{enumerate}

\PSsubsection{When to Use}

\PSyes{} Thermal noise in circuits\PSbr{}\PSyes{} Measurement errors\PSbr{}\PSyes{} Any sum of many independent random effects (CLT)\PSbr{}\PSyes{} Manufacturing variations\PSbr{}\PSyes{} When data histogram is bell-shaped and symmetric

\PSsubsection{When NOT to Use}

\PSno{} Strictly positive quantities (Gaussian allows negative values)\PSbr{}\PSno{} Heavily skewed data\PSbr{}\PSno{} Bounded quantities (Gaussian has infinite support)\PSbr{}\PSno{} Discrete counts\PSbr{}\PSno{} Heavy-tailed phenomena (use t-distribution or Cauchy)

\PSsection{3.10}{Gamma Distribution}

\PSsubsection{Physical Interpretation}

Models the \textbf{total waiting time until the $k-\mathrm{th}$ event} in a Poisson process, or any sum of $k$ independent exponential random variables.

\PSsubsection{Engineering Context}

\begin{itemize}
\tightlist
\item
  Total repair time for $k$ components (each with exponential repair time)
\item
  Time until $k-\mathrm{th}$ packet arrival
\item
  Aggregate service time for $k$ queued jobs
\item
  Rainfall amount modeling
\end{itemize}

\begin{PSbox}[unbreakable]{def}{Mathematical Definition}

$X \sim \mathrm{Gamma}(\alpha,\, \beta)$ where $\alpha$ = shape parameter, $\beta$ = scale parameter (or rate $\lambda = \frac{1}{\beta}$).

\end{PSbox}

\PSsubsection{PDF}

\PSdisplay{f_X(x) = \frac{x^{\alpha-1} e^{-x/\beta}}{\beta^\alpha \Gamma(\alpha)}, \quad x \geq 0}

where $\Gamma (\alpha) = (\alpha -1)$! for integer $\alpha$, and $\Gamma (\alpha) = \int_{0}^{\infty} t^{\alpha -1}e^{-t}dt$ in general.

\PSsubsection{CDF}

No closed form for general $\alpha$. Available via incomplete gamma function.

\PSsubsection{Mean}

\PSdisplay{E[X] = \alpha\beta}

\PSsubsection{Variance}

\PSdisplay{\text{Var}(X) = \alpha\beta^2}

\PSsubsection{Special Cases}

{\def\LTcaptype{none} % do not increment counter
\begin{longtable}[c]{@{}cc@{}}
\toprule\noalign{}
\rowcolor{PSnavy}\PSth{Parameters} & \PSth{Distribution} \\
\midrule\noalign{}
\endhead
\bottomrule\noalign{}
\endlastfoot
$\alpha = 1$ & $\mathrm{Exponential}(\beta)$ \\
$\alpha = \frac{n}{2},\, \beta = 2$ & Chi-square(n) \\
$\alpha$ = integer $k$ & $\mathrm{Erlang}(k,\, \beta)$ \\
\end{longtable}
}

\PSsubsection{When to Use}

\PSyes{} Sum of independent exponential lifetimes\PSbr{}\PSyes{} Waiting time for $k-\mathrm{th}$ Poisson event\PSbr{}\PSyes{} Skewed, non-negative data\PSbr{}\PSyes{} Flexible shape (adjustable skewness)

\PSsubsection{When NOT to Use}

\PSno{} Data can be negative\PSbr{}\PSno{} Data is symmetric (use Gaussian)\PSbr{}\PSno{} Simpler distribution fits (exponential for $k = 1$)

\PSsection{3.11}{Rayleigh Distribution}

\PSsubsection{Physical Interpretation}

Models the \textbf{magnitude (envelope)} of a $2D$ Gaussian random vector. When $X$ and $Y$ are independent $N(0,\, \sigma^{2})$, then $R = \sqrt{X^{2} + Y^{2}}$ follows a Rayleigh distribution.

\PSsubsection{Engineering Context}

\begin{itemize}
\tightlist
\item
  Envelope of narrowband noise (in-phase + quadrature)
\item
  Magnitude of wireless fading channel coefficient (Rayleigh fading)
\item
  Wind speed modeling
\item
  Distance error in $2D$ positioning (when errors are Gaussian in $x$ and $y$)
\end{itemize}

\begin{PSbox}[unbreakable]{def}{Mathematical Definition}

$R \sim \mathrm{Rayleigh}(\sigma)$, where $\sigma$ is the parameter (NOT the standard deviation of $R$).

\end{PSbox}

\PSsubsection{PDF}

\PSdisplay{f_R(r) = \frac{r}{\sigma^2} e^{-r^2/(2\sigma^2)}, \quad r \geq 0}

\PSsubsection{CDF}

\PSdisplay{F_R(r) = 1 - e^{-r^2/(2\sigma^2)}, \quad r \geq 0}

\PSsubsection{Mean}

\PSdisplay{E[R] = \sigma\sqrt{\frac{\pi}{2}} \approx 1.2533\sigma}

\PSsubsection{Variance}

\PSdisplay{\text{Var}(R) = \frac{4-\pi}{2}\sigma^2 \approx 0.4292\sigma^2}

\PSsubsection{Relationship to Other Distributions}

\begin{itemize}
\tightlist
\item
  $R^{2} \sim \mathrm{Exponential}(2\sigma^{2})$ — the power is exponentially distributed
\item
  If $X,\, Y \sim N(0,\, \sigma^{2})$ independent, then $\sqrt{X^{2}+Y^{2}} \sim \mathrm{Rayleigh}(\sigma)$
\end{itemize}

\PSsubsection{When to Use}

\PSyes{} Envelope/magnitude of complex Gaussian signal\PSbr{}\PSyes{} Rayleigh fading channel modeling (no line-of-sight)\PSbr{}\PSyes{} $2D$ distance with Gaussian errors in each dimension\PSbr{}\PSyes{} RSS (root-sum-square) of two independent Gaussian components

\PSsubsection{When NOT to Use}

\PSno{} Strong line-of-sight component present (use Rician)\PSbr{}\PSno{} Not modeling a magnitude/envelope\PSbr{}\PSno{} Components are not Gaussian or not equal-variance

\PSsection{3.12}{Chi-Square Distribution}

\PSsubsection{Physical Interpretation}

Models the \textbf{sum of squares of independent standard normal} random variables. Arises naturally in variance estimation and goodness-of-fit testing.

\PSsubsection{Engineering Context}

\begin{itemize}
\tightlist
\item
  Sample variance distribution (for Gaussian data)
\item
  Power of Gaussian noise (sum of squared components)
\item
  Goodness-of-fit test statistic
\item
  Confidence intervals for variance
\item
  Energy detection in spectrum sensing
\end{itemize}

\begin{PSbox}[unbreakable]{def}{Mathematical Definition}

If $Z_{1},\, Z_{2},\, \ldots,\, Z_{n}$ are independent $N(0,\,1)$, then:\PSdisplay{\chi^2 = Z_1^2 + Z_2^2 + \cdots + Z_n^2 \sim \chi^2(n)}

$n$ = degrees of freedom.

\end{PSbox}

\PSsubsection{PDF}

\PSdisplay{f_X(x) = \frac{x^{n/2-1} e^{-x/2}}{2^{n/2}\Gamma(n/2)}, \quad x \geq 0}

\PSsubsection{CDF}

No closed form for general $n$. Computed via incomplete gamma function.

\PSsubsection{Mean}

\PSdisplay{E[X] = n}

\PSsubsection{Variance}

\PSdisplay{\text{Var}(X) = 2n}

\PSsubsection{Relationship to Other Distributions}

\begin{itemize}
\tightlist
\item
  Chi-square(n) $= \mathrm{Gamma}\left(\frac{n}{2},\, 2\right)$
\item
  Chi-square(1) = square of $N(0,\,1)$
\item
  Chi-square(2) $= \mathrm{Exponential}\left(\frac{1}{2}\right)$
\item
  For large $n$: $\chi^{2}(n) \approx N(n,\, 2n)$ by CLT
\end{itemize}

\PSsubsection{When to Use}

\PSyes{} Sum of squared Gaussian random variables\PSbr{}\PSyes{} Variance testing and confidence intervals\PSbr{}\PSyes{} Goodness-of-fit tests\PSbr{}\PSyes{} Energy detection in $N$ dimensions

\PSsubsection{When NOT to Use}

\PSno{} Variables are not Gaussian\PSbr{}\PSno{} Variables are not independent\PSbr{}\PSno{} Variables are not standard normal (need scaling)

\PSsection{3.13}{Distribution Selection Guide}

\PSsubsection{Decision Tree: Discrete Distributions}

\PSfigure{m03-tree-discrete}

\PSsubsection{Decision Tree: Continuous Distributions}

\PSfigure{m03-tree-continuous}

\PSsubsection{Quick Recognition Patterns}

{\def\LTcaptype{none} % do not increment counter
\begin{longtable}[c]{@{}cc@{}}
\toprule\noalign{}
\rowcolor{PSnavy}\PSth{If you see…} & \PSth{Think…} \\
\midrule\noalign{}
\endhead
\bottomrule\noalign{}
\endlastfoot
Binary outcome & Bernoulli \\
``$n$ trials, $k$ successes'' & Binomial \\
``How many until first…'' & Geometric \\
``Events per interval'' & Poisson \\
``Equal chance anywhere in $[a,\,b]$'' & Uniform \\
``Time until next event'' & Exponential \\
``Sum of many effects'' or ``noise'' & Gaussian \\
``Total waiting time for $k$ events'' & Gamma \\
``Signal envelope'' or ``fading'' & Rayleigh \\
``Sum of squared normals'' & Chi-square \\
\end{longtable}
}

\PSsection{3.14}{Summary Table}

{\def\LTcaptype{none} % do not increment counter
\begin{longtable}[c]{@{}
  >{\raggedright\arraybackslash}p{0.20\linewidth}
  >{\centering\arraybackslash}p{0.12\linewidth}
  >{\centering\arraybackslash}p{0.21\linewidth}
  >{\centering\arraybackslash}p{0.17\linewidth}
  >{\centering\arraybackslash}p{0.20\linewidth}@{}}
\toprule\noalign{}
\rowcolor{PSnavy}\PSth{Distribution} & \PSth{Type} & \PSth{Parameters} & \PSth{Mean} & \PSth{Variance} \\
\midrule\noalign{}
\endhead
\bottomrule\noalign{}
\endlastfoot
Bernoulli & Discrete & $p$ & $p$ & $p(1-p)$ \\
Binomial & Discrete & $n,\, p$ & $np$ & $np(1-p)$ \\
Geometric & Discrete & $p$ & $\frac{1}{p}$ & $\frac{1-p}{p^{2}}$ \\
Neg. Binomial & Discrete & $r,\, p$ & $\frac{r}{p}$ & $\frac{r(1-p)}{p^{2}}$ \\
Poisson & Discrete & $\lambda$ & $\lambda$ & $\lambda$ \\
Uniform & Continuous & a, $b$ & $\frac{a+b}{2}$ & $\frac{(b-a)^{2}}{12}$ \\
Exponential & Continuous & $\lambda$ (or $\beta = \frac{1}{\lambda}$) & $\frac{1}{\lambda}$ & $\frac{1}{\lambda^{2}}$ \\
Gaussian & Continuous & $\mu,\, \sigma^{2}$ & $\mu$ & $\sigma^{2}$ \\
Gamma & Continuous & $\alpha,\, \beta$ & $\alpha \beta$ & $\alpha \beta^{2}$ \\
Rayleigh & Continuous & $\sigma$ & $\sigma \sqrt{\frac{\pi}{2}}$ & $\frac{(4-\pi)\sigma^{2}}{2}$ \\
Chi-square & Continuous & $n$ & $n$ & $2n$ \\
\end{longtable}
}

\PSsection{3.15}{Practice Problems}

\begin{PSbox}[unbreakable]{prob}{Problem 1: Distribution Identification}

Identify the appropriate distribution for each scenario:

\begin{enumerate}
\def\labelenumi{(\alph{enumi})}
\tightlist
\item
  Number of dropped calls in a 10-minute window (average: 2 per 10 min)
\item
  Voltage of thermal noise across a resistor
\item
  Time until a server crashes (constant failure rate)
\item
  Number of defective ICs out of a batch of 100 (defect rate 3\%)
\item
  Envelope of received signal through multipath with no LOS
\end{enumerate}

\PSsolution{} (a) $\mathrm{Poisson}(\lambda = 2)$, (b) $\mathrm{Gaussian}(0,\, \sigma^{2})$, (c) $\mathrm{Exponential}(\lambda)$, (d) $\mathrm{Binomial}(100,\, 0.03)$, (e) $\mathrm{Rayleigh}(\sigma)$

\end{PSbox}

\begin{PSbox}[unbreakable]{prob}{Problem 2: Exponential vs Poisson}

Packets arrive at a router as a Poisson process with rate $\lambda = 5$ packets/second.

\begin{enumerate}
\def\labelenumi{(\alph{enumi})}
\tightlist
\item
  What is $P(\text{more than 8 packets in 1 second})$?
\item
  What is $P(\text{no packet arrives in the next 0.5 seconds})$?
\item
  What distribution describes the inter-arrival time?
\item
  What is the mean inter-arrival time?
\end{enumerate}

\PSsolution{} 

\begin{enumerate}
\def\labelenumi{(\alph{enumi})}
\tightlist
\item
  $X \sim \mathrm{Poisson}(5)$: $P(X > 8) = 1 - \texttt{poisscdf}(8,\, 5) = 0.0681$
\item
  $T \sim \mathrm{Exp}(5)$: $P(T > 0.5) = e^{-5 \times 0.5} = e^{-2.5} = 0.0821$
\item
  Exponential with rate $\lambda = 5$
\item
  Mean $= \frac{1}{\lambda} = 0.2$ seconds
\end{enumerate}

\end{PSbox}

\begin{PSbox}[unbreakable]{prob}{Problem 3: Gaussian Probability}

A resistor has nominal value 1000\ensuremath{\Omega} with manufacturing variation modeled as $N(1000,\, 25) (\sigma = 5\,\Omega)$.

\begin{enumerate}
\def\labelenumi{(\alph{enumi})}
\tightlist
\item
  What fraction of resistors are within $\pm 10\,\Omega$ of nominal (spec: $990-1010\,\Omega$)?
\item
  What fraction fail the $\pm 1\%$ tolerance test?
\item
  What tolerance range contains 99.7\% of resistors?
\end{enumerate}

\PSsolution{} 

\begin{enumerate}
\def\labelenumi{(\alph{enumi})}
\tightlist
\item
  $P(990 \le X \le 1010) = P(-2 \le Z \le 2) = 0.9545 \to 95.45\%$
\item
  $\pm 1\% = \pm 10\,\Omega$ \ensuremath{\to} same as (a), fail rate $= 1 - 0.9545 = 4.55\%$
\item
  $\pm 3\sigma = \pm 15\,\Omega$ \ensuremath{\to} range $[985,\, 1015]\Omega$
\end{enumerate}

\emph{Next Module: {Module 4 — Expectation and Moments}}

\end{PSbox}
